% ============================================================
% PRN232 - Section 1: Introduction (Vietnamese & English)
% ============================================================
\section{\bi{Mở đầu}{Introduction}}
\label{sec:intro}

\bi{
Trong những năm gần đây, nhu cầu theo dõi nội dung truyền hình trực tuyến và tương tác trực tiếp (LiveStream)
tại Việt Nam đã thay đổi rõ rệt. Người dùng không còn chỉ ngồi trước
tivi lúc 21 giờ; họ xem các chương trình theo dõi ngẫu nhiên, tham gia phòng phát trực tiếp với bình luận và biểu cảm thời gian thực, đồng thời muốn được gợi ý nội dung phù hợp
với sở thích. Ở phía bên kia, các đài truyền hình và nền tảng nội
dung số cần một hệ thống lịch phát và bàn điều khiển phát sóng (Live Studio) thống nhất, đáng tin cậy và có thể
mở rộng để phục vụ đồng thời nhiều nền tảng: web, ứng dụng di động và
các thiết bị phát trực tuyến.

Bài báo này trình bày \textbf{OmniCast} --- một hệ thống lập lịch, truyền phát trực tiếp (LiveStream) và
phân phối nội dung số đa nền tảng, được nghiên cứu và phát triển tại \textbf{FPTU Quy Nhon AI Campus --- FU Cahata} cho học phần
\textbf{PRN232 --- Building Cross-Platform Back-End Application} dưới sự hướng dẫn chuyên môn của giảng viên (\emph{Mentor}) \textbf{Thầy Huỳnh Vang}.
Hệ thống gồm một back-end NestJS hiệu năng cao phục vụ đồng thời một cổng
web Next.js 15 (33 route) và một ứng dụng di động Flutter (31 màn hình), với lõi nghiệp vụ tập trung vào
bốn năng lực chủ chốt: hệ sinh thái LiveStream \& Live Studio thời gian thực hai chiều, lịch phát điện tử (EPG) 24/7 không bao giờ để trống,
cơ chế phát hiện xung đột lịch phát nguyên tử, và trợ lý biên tập AI có kiểm
duyệt theo chuẩn đạo đức nội dung.
}{
In recent years, the way Vietnamese audiences consume television and interactive live broadcasts (LiveStream) has
changed dramatically. Viewers no longer sit in front of a screen at a
fixed broadcast slot; they browse on-demand, join interactive live rooms with real-time chat and floating reactions, and expect content recommendations tailored to
their interests. Broadcasters and digital-content providers, in turn,
need a single, reliable and scalable scheduling and Live Studio control system that serves
several delivery surfaces at once: the web, mobile applications and
live-streaming devices.

This paper presents \textbf{OmniCast} --- a multi-platform digital
broadcast scheduling, real-time LiveStream and content delivery system developed at \textbf{FPTU Quy Nhon AI Campus --- FU Cahata} for the course
\textbf{PRN232 --- Building Cross-Platform Back-End Application} under the academic mentorship of \textbf{Mr. Huỳnh Vang}. The
system comprises a high-performance NestJS back end that serves both a
Next.js 15 web portal (33 routes) and a Flutter mobile application (31 screens), with the core
business logic concentrated on four capabilities: a bi-directional real-time LiveStream and Live Studio ecosystem, an electronic
programme guide (EPG) that is never empty, an atomic broadcast-schedule
conflict-detection mechanism, and an AI editorial assistant with
content-compliance validation.
}

\subsection{\bi{Động cơ và vấn đề cần giải quyết}{Motivation and Problem Statement}}

\bi{
Bốn vấn đề kỹ thuật chi phối toàn bộ thiết kế của hệ thống được xác định
từ yêu cầu thực tế của dự án:
}{
Four technical problems dominate the system's design, each derived
from concrete project requirements:
}

\begin{enumerate}[leftmargin=*,itemsep=2pt]
  \item \bi{\textbf{Truyền phát trực tiếp đa nguồn và điều hành phòng Live thời gian thực.} Một nền tảng truyền hình hiện đại phải cho phép đạo diễn phát sóng từ thư viện VOD, tự ghi hình từ Camera/Màn hình qua WebRTC WHIP, hoặc lấy từ kho nội dung công khai đã xác thực; đồng thời cho phép chuyển đổi chương trình đang phát ngay lập tức (\emph{on-the-fly program switching}) mà không làm rớt kết nối của hàng nghìn khán giả đang trò chuyện, thả cảm xúc và tặng quà Super Chat trong phòng Live.}{
  \textbf{Multi-source LiveStream ingestion and real-time studio orchestration.} A modern broadcast platform must allow producers to go live from curated VOD libraries, device cameras/screens over WebRTC WHIP, or verified public catalogs, while supporting zero-downtime on-the-fly program switching without dropping viewers who are actively chatting, sending floating reactions and donating via Super Chat.}

  \item \bi{\textbf{Lịch phát thưa thớt.} Một kênh truyền hình thực tế
  chỉ có khoảng một vài chương trình gốc mỗi ngày. Nếu dựng lưới EPG 24
  giờ trực tiếp từ bảng dữ liệu, phần lớn lưới sẽ trống và bộ lọc theo
  ngày ở giao diện người dùng trở nên vô nghĩa. Hệ thống cần một cơ
  chế điền lịch tự động, luôn đảm bảo mỗi kênh trong mỗi ngày đều đầy
  đủ 24 giờ nội dung hợp lệ.}{
  \textbf{Sparse schedules.} A real broadcast channel originates only a
  handful of productions per day. Building a 24-hour EPG grid directly
  from the event table leaves most of the grid empty and renders the
  day filter in the user interface meaningless. The system therefore
  requires an automatic filler that guarantees every channel-day grid
  is fully populated with valid content.}

  \item \bi{\textbf{Xung đột lịch phát.} Nếu hai chương trình cùng một
  kênh có khoảng thời gian chồng lấn, người xem sẽ không biết chương
  trình nào được phát. Cơ chế kiểm tra phải phát hiện được cả ba dạng
  giao thoa (đầu, đuôi, bao trùm toàn bộ) và phải chạy trong một câu
  truy vấn cơ sở dữ liệu để tránh điều kiện tranh chấp.}{
  \textbf{Schedule conflicts.} If two programmes on the same channel
  overlap, viewers cannot tell which one is actually on air. The
  validation must detect all three overlap forms (head, tail, full
  containment) and must execute as a single database query to avoid
  race conditions.}

  \item \bi{\textbf{Kiểm duyệt nội dung không tự động.} Quy trình biên
  tập nội dung (mô tả, thẻ, phân loại) là thủ công và tốn thời gian.
  Việc tự động hoá bằng mô hình ngôn ngữ lớn giúp giảm tải, nhưng đòi
  hỏi cơ chế dự phòng khi không có khóa API, và đầu ra phải được kiểm
  tra theo schema trước khi lưu.}{
  \textbf{Manual content moderation.} The editorial workflow (writing
  descriptions, tagging, categorising) is manual and time-consuming.
  Automating it with a large language model reduces the load, but
  requires a fallback path when no API key is present, and the output
  must be schema-validated before it is persisted.}
\end{enumerate}

\subsection{\bi{Đóng góp của bài báo}{Contributions}}

\bi{
Bài báo đóng góp các kết quả kỹ thuật sau:
\begin{enumerate}[leftmargin=*,itemsep=2pt]
  \item Kiến trúc phân tầng \textbf{NestJS + Prisma + PostgreSQL} với 19
  module nghiệp vụ, hơn 140 endpoint REST và tài liệu OpenAPI tự sinh,
  kèm lớp bảo vệ RBAC bốn vai trò (\texttt{GUEST}, \texttt{VIEWER}, \texttt{STAFF}, \texttt{ADMIN}).
  \item Kiến trúc \textbf{LiveStream \& Live Studio thời gian thực} kết hợp MediaMTX (WHIP/WHEP/HLS), Cloudinary Direct Upload và cổng Socket.IO \texttt{/live} hỗ trợ 4 chế độ phát sóng Studio, chuyển chương trình tức thời, trò chuyện có trả lời phân nhánh, biểu cảm hạt nổi 0ms, quyên góp Super Chat 6 mệnh giá và thông báo ghim từ đạo diễn.
  \item Thuật toán phát hiện xung đột lịch phát bằng một truy vấn SQL
  thô, tính toán thời điểm kết thúc hiệu lực bằng
  \texttt{COALESCE(endedAt, scheduledAt + duration)}.
  \item Cơ chế điền lịch EPG 24/7 với lựa chọn ngẫu nhiên có hạt giống
  (seeded PRNG) theo cặp (ngày, kênh), bảo đảm cùng một ngày luôn
  cho cùng một lịch.
  \item Quy trình biên tập AI gọi mô hình một lần duy nhất
  (\texttt{gpt-4o}), kiểm tra schema bằng Zod, có cơ chế thử lại và
  đường dự phòng heuristic.
  \item Microservice ghi nhật ký kiểm toán giao tiếp qua gRPC (\texttt{audit\_logger.proto}),
  tách biệt khỏi tiến trình API chính.
  \item Cổng web Next.js 15 với 33 route (bao gồm phòng Live 7 chủ đề giao diện, bàn điều khiển Studio 4 tab và quản trị AI Curator) cùng quy trình CI/CD 5 job kiểm định 482/482 kiểm thử tự động đạt 100\%.
\end{enumerate}
}{
This paper makes the following technical contributions:
\begin{enumerate}[leftmargin=*,itemsep=2pt]
  \item A layered \textbf{NestJS + Prisma + PostgreSQL} architecture
  with 19 business modules, 140+ REST endpoints and auto-generated
  OpenAPI documentation, protected by a four-role RBAC layer (\texttt{GUEST}, \texttt{VIEWER}, \texttt{STAFF}, \texttt{ADMIN}).
  \item A \textbf{real-time LiveStream \& Live Studio architecture} combining MediaMTX (WHIP/WHEP/HLS), Cloudinary Direct Upload and the Socket.IO \texttt{/live} gateway to support 4 studio broadcast modes, zero-downtime program switching, threaded live chat, 0ms floating emoji reactions, 6-tier Super Chat donations and pinned broadcaster announcements.
  \item A broadcast-conflict detection algorithm implemented as a
  single raw SQL query that computes the effective end time with
  \texttt{COALESCE(endedAt, scheduledAt + duration)}.
  \item A 24/7 EPG filler with seeded pseudo-random selection keyed on
  the (date, channel) pair, guaranteeing that the same day always
  renders the same schedule.
  \item An AI editorial pipeline that issues a single consolidated model
  call (\texttt{gpt-4o}), validates the result against a Zod schema,
  and provides retry plus heuristic-fallback paths.
  \item An audit-logging microservice communicating over gRPC (\texttt{audit\_logger.proto}), deployed
  independently of the main API process.
  \item A Next.js 15 web portal with 33 routes (including a 7-theme LiveStream lounge, a 4-tab Live Studio console and an AI Curator suite) backed by a five-job GitHub Actions CI/CD pipeline with 482/482 passing automated tests (100\%).
\end{enumerate}
}

\subsection{\bi{Tổ chức nhóm nghiên cứu và Phân công nhiệm vụ}{Project Team Organization and Role Allocation}}

\bi{
Dự án OmniCast được thực hiện tại \textbf{FPTU Quy Nhon AI Campus --- FU Cahata} dưới sự hướng dẫn của giảng viên \textbf{Thầy Huỳnh Vang}. Bảng~\ref{tab:team_prn} trình bày phân công nhiệm vụ chuyên sâu của 4 thành viên trong nhóm.
}{
The OmniCast project was executed at \textbf{FPTU Quy Nhon AI Campus --- FU Cahata} under the supervision of Mentor \textbf{Mr. Huỳnh Vang}. Table~\ref{tab:team_prn} summarises the responsibilities of the four team members.
}

\begin{table}[t]
\caption{\bi{Danh sách 4 thành viên nhóm dự án OmniCast (FPTU Quy Nhon AI Campus --- FU Cahata, Mentor: Thầy Huỳnh Vang).}{The 4 project team members of OmniCast (FPTU Quy Nhon AI Campus --- FU Cahata, Mentor: Mr. Huỳnh Vang).}}
\label{tab:team_prn}
\centering
\scriptsize
\begin{tabular}{@{}p{2.3cm}p{1.5cm}p{4.3cm}@{}}
\toprule
\textbf{\bi{Họ và tên}{Full Name}} & \textbf{\bi{Vai trò / MSSV}{Role / ID}} & \textbf{\bi{Phụ trách kỹ thuật chính}{Primary Engineering Ownership}} \\
\midrule
\textbf{Nguyễn Bình An} (\texttt{TomOutfit}) & \bi{Nhóm trưởng}{Team Leader} & \bi{Kiến trúc lõi NestJS 19 module, Bảo mật JWT \& RBAC 4 vai trò, Thuật toán chống xung đột lịch EPG, Direct Stream Engine, CI/CD \& Triển khai}{NestJS 19-module core architecture, JWT \& 4-role RBAC security, EPG conflict detection SQL, Direct Stream Engine, CI/CD \& Release} \\
\addlinespace[2pt]
\textbf{Trần Thanh Sơn} & \texttt{QE170080} & \bi{Kiến trúc vi dịch vụ gRPC Audit Logger (\texttt{audit\_logger.proto}), Cơ sở dữ liệu PostgreSQL Supabase (21 Prisma models), Telemetry \& Seed}{Standalone gRPC Audit Logger microservice (\texttt{audit\_logger.proto}), Supabase PostgreSQL (21 Prisma models), Telemetry \& Seeding} \\
\addlinespace[2pt]
\textbf{Nguyễn Vĩnh Phú} & \texttt{QE180136} & \bi{Trưởng nhóm Mobile Flutter (31 màn hình, 12 BLoC), Trình phát \texttt{OmniPlayer}, \texttt{MobileStudioRecorderScreen} (WebRTC WHIP), Offline SQLite}{Mobile Flutter Lead (31 screens, 12 BLoCs), \texttt{OmniPlayer} engine, \texttt{MobileStudioRecorderScreen} (WebRTC WHIP), Offline SQLite} \\
\addlinespace[2pt]
\textbf{Đồng Hoàng Nguyên} (\texttt{teng122}) & Full-Stack \& AI & \bi{Cổng Web Next.js 15 Cinema Dark (33 route), Cổng Socket.IO \texttt{/live}, Phòng Live 7 chủ đề, Bàn điều khiển Live Studio \& Module AI Curator}{Next.js 15 Cinema Dark Web Portal (33 routes), Socket.IO \texttt{/live} gateway, 7-theme Live Room, Live Studio Console \& AI Curator module} \\
\bottomrule
\end{tabular}
\end{table}

\subsection{\bi{Phạm vi bài báo}{Scope}}

\bi{
Bài báo này tập trung vào kiến trúc hệ thống, phần back-end NestJS, vi dịch vụ gRPC, cơ sở dữ liệu và cổng web Next.js 15 của hệ thống (thuộc học phần PRN232). Phần ứng dụng di động Flutter 31 màn hình được trình bày
đầy đủ trong báo cáo song hành thuộc học phần PRM393.}{
This paper concentrates on the overall system architecture, the NestJS back end, the gRPC microservice, the database schema and the Next.js 15 web portal (PRN232 deliverables). The 31-screen Flutter mobile application is
presented in full detail in the companion PRM393 report.
}
